\scanpage{307}
\mspage{257}

Donc le cor.~1 résulte \unsure{de} la formule
\[
p_{\widetilde\Psi*}=\varepsilon_3^*,
\]
qui \deleted{résulte} elle-même est un cas particulier du

\uline{Lemme} Soit $I$ dans $\mathrm{Cat}$, avec objet final $e$, considérons l'inclusion
\[
\varepsilon\colon e\longrightarrow I,\quad\text{et }p\colon I\mathrel{\unsure{-|\to}}e
\]

Soit $(M,W)$ une catégorie de modèles, et considérons
\[
\deleted{\unsure{M}}\ \mathrm{Ho}_W(I)\xleftarrow{p^*}\mathrm{Ho}_W .
\]
Ce foncteur admet un adjoint à gauche $p_*$, lequel est isom. à $\varepsilon^*$.

\uline{Dém} Considérons le couple de foncteurs adjoints
\diagram{\begin{tikzcd}[ampersand replacement=\&]
M(I)\arrow[r,shift right,"{p_*=\varprojlim_I}"']\&M\arrow[l,shift right,"{p^*}"']
\end{tikzcd}\ ;}
comme $I$ admet \deleted{\unsure{de}} l'objet final $e$, le foncteur $p_*$ est isomorphe à $\varepsilon^*$. Ces foncteurs $p^*,\varepsilon^*$ sont compatibles aux localiseurs\unsure{,|} $W(I),W$, donc on conclut par le lemme page 248.

Il ne reste plus qu'à prouver les lemmes 1 et 2. Je le ferai demain.
% archive sheet number: 307
